% Generated draft: requires mathematical and source review.
\documentclass[11pt]{article}
\usepackage[margin=25mm]{geometry}
\usepackage{fontspec}
\setmainfont{Latin Modern Roman}
\usepackage{amsmath,amssymb,mathtools,unicode-math}
\setmathfont{Latin Modern Math}
\usepackage{enumitem,needspace,xurl,hyperref}
\setlength{\parindent}{0pt}
\setlength{\parskip}{5pt}
\setlength{\emergencystretch}{3em}
\newcommand{\N}{\mathbb{N}}
\newcommand{\Z}{\mathbb{Z}}
\newcommand{\Q}{\mathbb{Q}}
\newcommand{\R}{\mathbb{R}}
\newcommand{\C}{\mathbb{C}}
\newcommand{\F}{\mathbb{F}}
\newcommand{\A}{\mathbb{A}}
\newcommand{\PP}{\mathbb{P}}
\newcommand{\E}{\mathbb{E}}
\newcommand{\eps}{\varepsilon}
\newcommand{\dd}{\,\mathrm{d}}
\begin{document}
\section*{UnsolvedMath ID: 30001521}
\textbf{Source problem code:} catalogue entry
\subsection{Definitions and mathematical statement}
Let \(X\) be a finite connected complex with basepoint \(x_0\). Assume that \(\pi_1(X,x_0)\) is abelian and that its natural action on \(\pi_n(X,x_0)\) is trivial for every \(n\geq 2\). Following the finite-dimensional wording of the cleaned source record, assume
\[
\dim_{\Q}\left(\bigoplus_{n\geq 1}\pi_n(X,x_0)\otimes_{\Z}\Q\right)<\infty.
\]
The term in degree one is defined because the fundamental group is abelian. The original wording, ``finitely generated rational homotopy,'' is not explicitly explained in the supplied record; whether it means this total finite-dimensional condition requires source review.

For a prime \(p\) and an integer \(r\geq 0\), write \(T^r=(S^1)^r\) and \(E_{p,r}=(\Z/p\Z)^r\), with the usual torus topology and discrete topology, respectively. Interpret an action here as a continuous action on \(X\) itself. An action of a group \(G\) is free if \(g\cdot x=x\) implies \(g=e\) for every \(g\in G\) and \(x\in X\). Define, using the maxima specified in the source,
\[
\operatorname{rk}_T(X)=\max\{r\in\Z_{\geq 0}:T^r\text{ acts freely on }X\},
\qquad
\operatorname{rk}_p(X)=\max\{r\in\Z_{\geq 0}:E_{p,r}\text{ acts freely on }X\}.
\]
For \(r=0\), the acting group is trivial.

\textbf{Comparison problem.} Determine the relationship between \(\operatorname{rk}_T(X)\) and \(\operatorname{rk}_p(X)\) as \(X\) ranges over complexes satisfying these hypotheses and \(p\) ranges over primes, including any dependence on \(p\).

The supplied source poses a comparison question without specifying a single conjectured equality or inequality. Its dated curation describes an asymptotic upper bound for free \(p\)-rank at sufficiently large primes and counterexamples to a universal reverse inequality, but supplies neither the bound's formula nor the reverse inequality's precise quantifiers. These descriptions therefore cannot be promoted here to precise mathematical assertions. The category of complex and any additional restrictions on admissible actions also require confirmation from the source.
\subsection{Short English statement}
For a finite connected complex with abelian fundamental group acting trivially on higher homotopy and finite-dimensional total rational homotopy, compare the largest ranks of a torus and an elementary abelian \(p\)-group that can act freely on the complex. The supplied record leaves the precise comparison target unspecified and labels the problem partially solved; that status is supplied curation, not an independently verified conclusion.
\subsection{Sources}
\begin{itemize}
\item Supplied local UnsolvedMath registry record, permanent ID \(30001521\), title \emph{Free Toral Rank versus Free \(p\)-Rank}, source problem code \texttt{catalogue entry}, in the Oberwolfach Reports Open Problems collection. The record gives proposed year 2010 and names no individual proposer.
\item Supplied report citation: \emph{Cohomology of Finite Groups: Interactions and Applications} (2011), Oberwolfach Reports, pp.\ 1885--1952, DOI \texttt{10.4171/owr/2010/32}. This is also the supplied literature reference labeled ``Original Oberwolfach problem report'': \url{https://doi.org/10.4171/owr/2010/32}.
\item Supplied literature reference: Hanke, \emph{The stable free rank of symmetry of products of spheres}, \url{https://doi.org/10.1007/s00222-009-0197-3}.
\item Literature assessment dated exactly \texttt{2026-08-21}: the supplied curation labels the problem partially solved, attributes an asymptotic upper bound for free \(p\)-rank at all sufficiently large primes to the report, records counterexamples to a universal reverse inequality, and says that the full relationship remains open. This is supplied dated curation, not independently verified; the background expressly requests independent verification of time-sensitive status and cited claims.
\item The registry labels its statement \texttt{corrected\_verified} and reports that it was checked against the official Oberwolfach report. This verification claim is supplied metadata; no independent source consultation is asserted here.
\end{itemize}
\end{document}
